\rhead{ML601: Computer Vision}
\phantomsection 
\section*{Computer Vision (ML601)}
\label{major:ML_6}

\noindent \small{\hyperref[main:scheme]{$\leftarrow$ Back to Scheme}} \vspace{0.5cm}

\noindent \textbf{Credits:} 3 (2-0-2)

\subsection*{Theory}
\noindent\textbf{Unit 1} \\
\textit{Digital Image Fundamentals and Processing:} Image formation and digitization (sampling, quantization), Pixel representations and color spaces (RGB, HSV, Lab, grayscale), Spatial domain filtering (linear, nonlinear filters, smoothing, sharpening), Edge detection (Sobel, Prewitt, Canny, Laplacian of Gaussian), Histogram equalization and contrast enhancement, Morphological operations (erosion, dilation, opening, closing).

\vspace{0.3cm}

\noindent\textbf{Unit 2} \\
\textit{Frequency Domain Processing and Transforms:} 2D Discrete Fourier Transform (DFT) and properties, Fast Fourier Transform (FFT) implementation, High-pass/low-pass filtering in frequency domain, Homomorphic filtering for illumination correction, Discrete Cosine Transform (DCT) for compression, Wavelet transforms (Haar, Daubechies), Multi-resolution analysis and pyramid representations.

\vspace{0.3cm}

\noindent\textbf{Unit 3} \\
\textit{Feature Detection, Extraction, and Description:} Corner detection (Harris, Shi-Tomasi), Blob detection (LoG, DoG, Hessian), SIFT (scale-space extrema, keypoint description), SURF and ORB for real-time applications, HOG (Histogram of Oriented Gradients) for pedestrian detection, Feature matching (nearest neighbor, FLANN), RANSAC for robust estimation, Bag-of-visual-words model.

\vspace{0.3cm}

\noindent\textbf{Unit 4} \\
\textit{Camera Geometry and 3D Vision:} Pinhole camera model and intrinsic/extrinsic parameters, Camera calibration (Zhang's method, checkerboard patterns), Epipolar geometry and fundamental matrix, Stereo vision (disparity maps, rectification), Structure from Motion (SfM) pipeline, SLAM fundamentals (visual odometry, bundle adjustment), Depth estimation from monocular cues.

\vspace{0.3cm}

\noindent\textbf{Unit 5} \\
\textit{Deep Learning for Computer Vision:} CNN architectures for vision (AlexNet, VGG, ResNet, EfficientNet), Object detection (R-CNN family, YOLO, SSD, RetinaNet), Semantic segmentation (FCN, U-Net, DeepLab), Instance segmentation (Mask R-CNN), Visual transformers (ViT, Swin Transformer), Self-supervised learning (SimCLR, DINO), 3D vision with PointNet/PointNet++, Video analysis (optical flow, 3D CNNs).

\newpage

\subsection*{Practical}
\noindent\textbf{Lab 1: Image Processing Pipeline} \\
Focuses on spatial/frequency domain filtering, histogram equalization, and morphological operations.

\vspace{0.2cm}

\noindent\textbf{Lab 2: Edge Detection and Feature Extraction} \\
Focuses on Canny edge detection, HOG descriptors, and contour analysis.

\vspace{0.2cm}

\noindent\textbf{Lab 3: SIFT/SURF Feature Matching} \\
Focuses on keypoint detection, descriptor extraction, and robust matching with RANSAC.

\vspace{0.2cm}

\noindent\textbf{Lab 4: Camera Calibration} \\
Focuses on checkerboard calibration, intrinsic/extrinsic parameter estimation, and distortion correction.

\vspace{0.2cm}

\noindent\textbf{Lab 5: Stereo Vision and Depth Maps} \\
Focuses on stereo rectification, disparity computation, and 3D point cloud generation.

\vspace{0.2cm}

\noindent\textbf{Lab 6: Object Detection with YOLO} \\
Focuses on YOLO model training, custom dataset annotation, and real-time inference.

\vspace{0.2cm}

\noindent\textbf{Lab 7: Semantic Segmentation} \\
Focuses on U-Net/DeepLab implementation for pixel-wise classification.

\vspace{0.2cm}

\noindent\textbf{Lab 8: Visual SLAM Implementation} \\
Focuses on ORB-SLAM or DSO for real-time visual odometry and mapping.

\vspace{0.2cm}

\noindent\textbf{Lab 9: 3D Vision with Point Clouds} \\
Focuses on PointNet for 3D object classification and segmentation.

\vspace{0.2cm}

\noindent\textbf{Lab 10: Complete CV Robotics Application} \\
Focuses on end-to-end vision system integrating detection, tracking, and 3D perception.

\vspace{0.2cm}

\newpage
\rhead{Curriculum Schema}
